% AA-58
\chapter{Półproste algebry Liego. Forma Killinga i kryteria Cartana}
\label{aa:ch:semisimple-lie}

W rozdziale~\ref{aa:ch:lie-algebras} rozróżniliśmy
algebry nilpotentne i rozwiązalne. Teraz wyodrębnimy
największy ideał rozwiązalny i zbadamy algebry,
w których ideał ten jest zerowy. Kluczowym narzędziem
będzie forma dwuliniowa zbudowana ze śladów operatorów
reprezentacji dołączonej. W całym rozdziale algebry
Liego są skończenie wymiarowe.

\section{Radykał rozwiązalny}

\begin{lemat}[Rozwiązalność rozszerzenia]
\label{aa:lem:solvable-extension}
Niech $\mathfrak g$ będzie algebrą Liego nad ciałem $K$,
a $\mathfrak a$ jej ideałem. Jeśli $\mathfrak a$ oraz
$\mathfrak g/\mathfrak a$ są rozwiązalne, to
$\mathfrak g$ jest rozwiązalna.
\end{lemat}
\begin{proof}
Niech $q:\mathfrak g\to\mathfrak g/\mathfrak a$
będzie rzutem ilorazowym. Z rozwiązalności ilorazu
istnieje $r\ge0$, dla którego jego $r$-ty wyraz
szeregu pochodnego jest zerowy. Na mocy
Stwierdzenia~\ref{aa:prop:lie-series-images}
$q(\mathfrak g^{(r)})=(\mathfrak g/\mathfrak a)^{(r)}=0$,
czyli $\mathfrak g^{(r)}\subset\mathfrak a$.
Indukcja po $s\ge0$ daje
$\mathfrak g^{(r+s)}\subset\mathfrak a^{(s)}$:
dla $s=0$ już to wykazano, a przejście do $s+1$
wynika z wzięcia nawiasu obu stron z samymi sobą.
Ponieważ pewne $\mathfrak a^{(s)}=0$, otrzymujemy
$\mathfrak g^{(r+s)}=0$.
\end{proof}

\begin{twierdzenie}[Istnienie radykału rozwiązalnego]
\label{aa:tw:lie-radical}
Niech $\mathfrak g$ będzie skończenie wymiarową algebrą
Liego nad ciałem $K$. Istnieje dokładnie jeden
największy (względem inkluzji) ideał rozwiązalny
$\operatorname{rad}\mathfrak g$.
Każdy automorfizm $\mathfrak g$ zachowuje ten ideał.
Iloraz $\mathfrak g/\operatorname{rad}\mathfrak g$
nie ma niezerowych ideałów rozwiązalnych.
\end{twierdzenie}
\begin{proof}
Niech $\mathfrak a,\mathfrak b$ będą ideałami
rozwiązalnymi. Ich suma algebraiczna jest ideałem,
bo $[x,a+b]=[x,a]+[x,b]\in\mathfrak a+\mathfrak b$.
Odwzorowanie
$\mathfrak b\to(\mathfrak a+\mathfrak b)/\mathfrak a$,
$b\mapsto b+\mathfrak a$, jest surjektywnym
homomorfizmem: klasa $a+b$ jest klasą $b$.
Iloraz jest rozwiązalny na mocy
Stwierdzenia~\ref{aa:prop:lie-series-images}.
Lemat~\ref{aa:lem:solvable-extension} zastosowany
do $\mathfrak a+\mathfrak b$ i jej ideału
$\mathfrak a$ dowodzi rozwiązalności sumy.

Wybieramy ideał rozwiązalny $\mathfrak r$ o największym
wymiarze; istnieje, ponieważ ideał zerowy jest
rozwiązalny, a zbiór możliwych wymiarów jest skończony.
Dla dowolnego ideału rozwiązalnego $\mathfrak b$
suma $\mathfrak r+\mathfrak b$ jest rozwiązalna,
więc jej wymiar nie przekracza $\dim\mathfrak r$.
Z inkluzji $\mathfrak r\subset\mathfrak r+\mathfrak b$
wynika równość, a stąd $\mathfrak b\subset\mathfrak r$.
To dowodzi istnienia i jednoznaczności.
Jeśli $f$ jest automorfizmem, $f(\mathfrak r)$
jest ideałem (dla $x=f(y)$ i $a\in\mathfrak r$
mamy $[x,f(a)]=f([y,a])$) oraz jest rozwiązalny
jako obraz $\mathfrak r$. Z maksymalności
$f(\mathfrak r)\subset\mathfrak r$ i z równości
wymiarów wynika równość.

Niech $\mathfrak c$ będzie ideałem rozwiązalnym
w $\mathfrak g/\mathfrak r$. Jego przeciwobraz
$\mathfrak d$ przez rzut $q$ jest ideałem, bo
$q([x,d])=[q(x),q(d)]\in\mathfrak c$. Zawiera
$\mathfrak r$, a $\mathfrak d/\mathfrak r\simeq
\mathfrak c$ na mocy twierdzenia o izomorfizmie
(Twierdzenie~\ref{aa:tw:lie-quotient}).
Lemat~\ref{aa:lem:solvable-extension} daje rozwiązalność
$\mathfrak d$, stąd $\mathfrak d\subset\mathfrak r$,
czyli $\mathfrak c=0$.
\end{proof}

\begin{definicja}[Prostota i półprostota]
\label{aa:def:semisimple-lie}
Algebra Liego $\mathfrak g$ nad ciałem $K$ jest prosta,
jeśli jest nieabelowa i jej jedynymi ideałami są
$0$ oraz $\mathfrak g$. Skończenie wymiarowa algebra
jest półprosta, jeśli $\operatorname{rad}\mathfrak g=0$.
\end{definicja}

Algebra prosta jest półprosta: gdyby była rozwiązalna,
jej niezerowa algebra pochodna, będąca ideałem,
musiałaby być równa całej algebrze, co przeczy
rozwiązalności. Algebra zerowa jest półprosta,
lecz nie prosta. Każda półprosta algebra ma zerowe
centrum, ponieważ centrum jest abelowym ideałem
(na mocy Stwierdzenia~\ref{aa:prop:adjoint-lie} jest
jądrem reprezentacji dołączonej).

\begin{przyklad}[Radykał w przykładach z AA-57]
\label{aa:prz:radical-examples}
Nad ciałem charakterystyki zero algebra Heisenberga
i algebra $\mathfrak b_n$ macierzy górnotrójkątnych
są rozwiązalne, więc ich radykały są całymi algebrami
(Przykłady~\ref{aa:prz:triangular-lie}
i~\ref{aa:prz:heisenberg-sl2}).
Algebra $\mathfrak{sl}_2$ jest prosta. Istotnie,
w bazie $h,e,f$ z Przykładu~\ref{aa:prz:heisenberg-sl2}
operator $\operatorname{ad}_h$ ma wartości własne
$0,2,-2$. Dla dowolnego niezerowego ideału
$\mathfrak a$ i niezerowego $u\in\mathfrak a$
każdy wielomian w $\operatorname{ad}_h$ zachowuje
$\mathfrak a$. Projekcje na trzy proste własne
są takimi wielomianami: dla wartości własnej $\lambda$
wystarczy użyć $\prod_{\mu\ne\lambda}
(t-\mu)/(\lambda-\mu)$, gdzie $\mu$ przebiega
pozostałe dwie wartości. Mianowniki są niezerowe
w charakterystyce zero. Co najmniej jedna projekcja
$u$ jest niezerowa, więc $\mathfrak a$ zawiera
co najmniej jeden z $h,e,f$. Jeśli zawiera $e$ lub $f$,
nawias z drugim daje $h$ lub $-h$. Jeśli zawiera $h$,
nawiasy z $e,f$ dają $2e,-2f$. W każdym przypadku
$\mathfrak a$ zawiera bazę, zatem $\mathfrak a=
\mathfrak{sl}_2$. Radykał $\mathfrak{sl}_2$ jest zerowy.
\end{przyklad}

\begin{center}
\begin{tikzpicture}[every node/.style={font=\small}]
 \node[draw,rounded corners,minimum width=32mm,minimum height=10mm] (g) at (0,0) {$\mathfrak g$};
 \node[draw,rounded corners,minimum width=32mm,minimum height=10mm] (r) at (-4,0) {$\operatorname{rad}\mathfrak g$};
 \node[draw,rounded corners,minimum width=32mm,minimum height=10mm] (q) at (4,0) {$\mathfrak g/\operatorname{rad}\mathfrak g$};
 \draw[-{Stealth}] (r)--node[above]{inkluzja} (g);
 \draw[-{Stealth}] (g)--node[above]{rzut} (q);
 \node at (-4,-.8) {rozwiązalna};
 \node at (4,-.8) {półprosta};
\end{tikzpicture}

\small Radykał rozwiązalny i iloraz półprosty.
Rysunek nie zakłada rozkładu $\mathfrak g$ na sumę ideałów.
\end{center}

\section{Formy niezmiennicze i forma Killinga}

\begin{definicja}[Forma niezmiennicza]
\label{aa:def:invariant-lie-form}
Niech $\mathfrak g$ będzie algebrą Liego nad $K$.
Symetryczna forma dwuliniowa $B:\mathfrak g\times
\mathfrak g\to K$ jest niezmiennicza, jeśli
$B([x,y],z)=B(x,[y,z])$ dla wszystkich $x,y,z$.
Jej jądro to $\{x:B(x,y)=0\ \forall y\in\mathfrak g\}$.
Forma jest niezdegenerowana, jeśli jądro jest zerowe.
Dla podprzestrzeni $U$ oznaczamy
$U^\perp=\{x:B(x,u)=0\ \forall u\in U\}$.
\end{definicja}

\begin{lemat}
\label{aa:lem:lie-orthogonal-ideals}
Niech $B$ będzie symetryczną niezmienniczą formą
dwuliniową na skończenie wymiarowej algebrze Liego
$\mathfrak g$. Jeśli $\mathfrak a$ jest ideałem,
to $\mathfrak a^\perp$ jest ideałem.
W szczególności jądro $B=\mathfrak g^\perp$ jest ideałem.
Jeśli $B$ jest niezdegenerowana, to
$[\mathfrak a,\mathfrak a^\perp]=0$.
\end{lemat}
\begin{proof}
Dla $x\in\mathfrak g$, $u\in\mathfrak a^\perp$,
$a\in\mathfrak a$ mamy
$B([x,u],a)=-B(u,[x,a])=0$; pierwsza równość
wynika z antysymetrii i niezmienniczości,
druga z $[x,a]\in\mathfrak a$.
Zatem $[x,u]\in\mathfrak a^\perp$.
Dla dowolnego $y\in\mathfrak g$ zachodzi
$B([u,a],y)=B(u,[a,y])=0$.
Jeśli forma jest niezdegenerowana, to $[u,a]=0$.
\end{proof}

\begin{definicja}[Forma śladu i forma Killinga]
\label{aa:def:killing}
Dla reprezentacji $\rho:\mathfrak g\to\mathfrak{gl}(V)$
w skończenie wymiarowej przestrzeni $V$ określamy
$B_\rho(x,y)=\operatorname{tr}_V(\rho(x)\rho(y))$.
Formą Killinga algebry $\mathfrak g$ nazywamy
formę śladu reprezentacji dołączonej:
\[
 \kappa_{\mathfrak g}(x,y)=
 \operatorname{tr}_{\mathfrak g}(\operatorname{ad}_x
 \operatorname{ad}_y).
\]
\end{definicja}

\begin{stwierdzenie}
\label{aa:prop:killing-properties}
Każda forma $B_\rho$ z Definicji~\ref{aa:def:killing}
jest symetryczna i niezmiennicza. Jeśli
$\mathfrak a\triangleleft\mathfrak g$ jest ideałem,
to $\kappa_{\mathfrak g}|_{\mathfrak a\times\mathfrak a}
=\kappa_{\mathfrak a}$.
\end{stwierdzenie}
\begin{proof}
Dwuliniowość wynika z liniowości $\rho$ i śladu.
Symetria wynika z $\operatorname{tr}(AB)=
\operatorname{tr}(BA)$. Dla $A=\rho(x),B=\rho(y),
C=\rho(z)$, korzystając z reprezentacji i cykliczności
śladu, otrzymujemy
\[
 \operatorname{tr}([A,B]C)=\operatorname{tr}(ABC-BAC)
 =\operatorname{tr}(A(BC-CB)).
\]
Cykliczność dla trzech czynników wynika z tej dla
dwóch po potraktowaniu iloczynu dwóch macierzy
jako jednego czynnika. To daje niezmienniczość.

Dla $x\in\mathfrak a$ operator $\operatorname{ad}_x$
na $\mathfrak g$ ma obraz w $\mathfrak a$, ponieważ
$\mathfrak a$ jest ideałem. W bazie powstałej przez
uzupełnienie bazy $\mathfrak a$ do bazy $\mathfrak g$
macierz ma zatem postać
$\begin{pmatrix}\operatorname{ad}^{\mathfrak a}_x&*\\0&0\end{pmatrix}$.
Analogiczny zapis dla $y\in\mathfrak a$ pokazuje,
że iloczyn ma na przekątnej bloki
$\operatorname{ad}^{\mathfrak a}_x
\operatorname{ad}^{\mathfrak a}_y$ oraz $0$.
Równość śladów daje tezę.
\end{proof}

\section{Kryterium rozwiązalności Cartana}

Najpierw dowiedziemy wersji macierzowej.
Istotną trudnością jest przejście od zanikania
śladów do nilpotentności operatorów. Nie wystarcza
samo $\operatorname{tr}(X^2)=0$ nad ciałem zespolonym.

\begin{twierdzenie}[Kryterium macierzowe Cartana]
\label{aa:tw:cartan-linear}
Niech $K$ będzie algebraicznie domkniętym ciałem
charakterystyki zero, $V$ skończenie wymiarową
przestrzenią wektorową nad $K$, a $\mathfrak l$
podalgebrą Liego $\mathfrak{gl}(V)$. Wówczas $\mathfrak l$ jest
rozwiązalna wtedy i tylko wtedy, gdy
\[
 \operatorname{tr}_V(XY)=0\qquad
 (X\in[\mathfrak l,\mathfrak l],\ Y\in\mathfrak l).
\]
\end{twierdzenie}
\begin{proof}
Jeśli $V=0$, algebra jest zerowa i obie własności
zachodzą. Załóżmy $V\ne0$. Jeśli $\mathfrak l$
jest rozwiązalna, Wniosek~\ref{aa:cor:lie-flag}
stosuje się do reprezentacji naturalnej: wszystkie
operatory z $\mathfrak l$ są w pewnej bazie
górnotrójkątne, a z jej algebry pochodnej ściśle
górnotrójkątne. Iloczyn macierzy ściśle
górnotrójkątnej z górnotrójkątną ma zerową
przekątną: w $(i,i)$-tym wpisie jedyny możliwy
indeks niezerowego składnika to $i$, lecz pierwszy
czynnik ma wpis $(i,i)$ równy zero. Ślad znika.

Odwrotnie załóżmy warunek śladowy i ustalmy
$X\in[\mathfrak l,\mathfrak l]$. Niech
$\lambda_1,\ldots,\lambda_s$ będą parami różnymi
wartościami własnymi $X$, a $V_i$ odpowiednimi
uogólnionymi przestrzeniami własnymi. Nad $K$
mamy $V=\bigoplus_{i=1}^s V_i$, a
$X|_{V_i}=\lambda_i I+N_i$, gdzie $N_i$ jest
nilpotentny (rozkład Jordana operatora).
Dokładniej, dla $d=\dim V$ możemy przyjąć
$V_i=\ker(X-\lambda_i I_V)^d$. Wielomian
charakterystyczny rozkłada się nad $K$ na czynniki
liniowe, więc rozkład na uogólnione przestrzenie
własne obejmuje całe $V$. Każde $V_i$ jest
niezerowe, gdyż zawiera wektor własny odpowiadający
$\lambda_i$, i $N_i^d=0$. Oznaczamy $m_i=\dim V_i>0$.
Niech $E=\operatorname{span}_{\mathbb Q}
\{\lambda_1,\ldots,\lambda_s\}\subset K$;
ciało $K$ zawiera kopię $\mathbb Q$ dzięki
charakterystyce zero. Dla dowolnego liniowego
nad $\mathbb Q$ funkcjonału $f:E\to\mathbb Q$
określamy $Y_f\in\operatorname{End}_K(V)$
przez $Y_f|_{V_i}=f(\lambda_i)I_{V_i}$.
Suma prosta zapewnia jednoznaczność takiego operatora
i jego liniowość nad $K$, mimo że sam funkcjonał
$f$ jest liniowy tylko nad $\mathbb Q$.
Nie zakładamy, że $Y_f\in\mathfrak l$.

Wykażemy jednak
$[Y_f,\mathfrak l]\subset[\mathfrak l,\mathfrak l]$.
Przestrzeń $\operatorname{End}_K(V)$ rozkłada się
na sumę prostą bloków $\operatorname{Hom}_K(V_j,V_i)$
(zapisanych jako operatory zerowe na pozostałych
składnikach). Na takim bloku
\[
 \operatorname{ad}_X(A)=(\lambda_i-\lambda_j)A
                         +N_iA-AN_j.
\]
Operatory $P(A)=N_i A$ i $Q(A)=AN_j$ na tym bloku
komutują, gdyż $P(Q(A))=N_iAN_j=Q(P(A))$.
Jeśli $N_i^{r_i}=0$ i $N_j^{r_j}=0$, to dla
$k=r_i+r_j-1$ mamy
\[
 (P-Q)^k(A)=\sum_{a=0}^k(-1)^a
      \binom{k}{a}N_i^{k-a}AN_j^a=0.
\]
Istotnie, gdy $a<r_j$, zachodzi $k-a\ge r_i$,
a gdy $a\ge r_j$, znika prawy czynnik. W szczególności dla
$d=\dim V$ potęga $2d-1$ zeruje część nilpotentną
na każdym bloku. Na tym samym bloku
$\operatorname{ad}_{Y_f}$ działa przez mnożenie
przez $f(\lambda_i)-f(\lambda_j)=f(\lambda_i-\lambda_j)$.

Niech $\Delta$ będzie zbiorem wszystkich różnic
$\lambda_i-\lambda_j$; jest podzbiorem $E$
i zawiera $0$. Jeśli ta sama różnica pojawia się
dla kilku par indeksów, przypisujemy jej tę samą
wartość $f(\delta)$, więc warunki interpolacyjne
są zgodne.
W pierścieniu $K[t]$ wielomiany $(t-\delta)^{2d}$
dla różnych $\delta\in\Delta$ są parami względnie
pierwsze: ich jedyne nierozkładalne czynniki są
różne. Chińskie twierdzenie o resztach daje
wielomian $p(t)$ spełniający
\[
 p(t)\equiv f(\delta)\pmod{(t-\delta)^{2d}}
 \qquad(\delta\in\Delta).
\]
Dla $\delta=0$ mamy $f(0)=0$, więc $p(0)=0$.
Dla każdej $\delta$ istnieje $q_\delta\in K[t]$
taki, że $p(t)-f(\delta)=(t-\delta)^{2d}q_\delta(t)$.
Podstawmy $T=\operatorname{ad}_X$ ograniczony do
bloku o różnicy $\delta$. Na tym bloku
$(T-\delta I)^{2d}=0$, a podstawienie do równości
wielomianów daje
$p(T)-f(\delta)I=(T-\delta I)^{2d}q_\delta(T)=0$.
Zatem $p(T)$ jest dokładnie mnożeniem przez
$f(\delta)$, także gdy część nilpotentna $T$ jest
niezerowa. Na wszystkich blokach otrzymujemy
$p(\operatorname{ad}_X)=\operatorname{ad}_{Y_f}$.
Ponieważ $X\in\mathfrak l$, operator
$\operatorname{ad}_X$ przeprowadza $\mathfrak l$
w $[\mathfrak l,\mathfrak l]$, a ten ideał w siebie.
Dla każdego $k\ge1$ indukcja daje
$(\operatorname{ad}_X)^k(\mathfrak l)
\subset[\mathfrak l,\mathfrak l]$:
pierwszy krok wynika z $X\in\mathfrak l$, następne
z własności ideału. Ponieważ $p(0)=0$, operator
$p(\operatorname{ad}_X)$ jest kombinacją liniową
tych dodatnich potęg. Stąd
$[Y_f,\mathfrak l]=p(\operatorname{ad}_X)(\mathfrak l)
\subset[\mathfrak l,\mathfrak l]$.

Z definicji algebry pochodnej istnieją
$A_k,B_k\in\mathfrak l$ takie, że
$X=\sum_{k=1}^r[A_k,B_k]$ (współczynniki liniowe
można włączyć do pierwszych czynników).
Cykliczność śladu daje dla każdego $k$
\[
\begin{aligned}
 \operatorname{tr}([A_k,B_k]Y_f)
 &=\operatorname{tr}(A_kB_kY_f-B_kA_kY_f)\\
 &=\operatorname{tr}(A_k(B_kY_f-Y_fB_k)).
\end{aligned}
\]
Mamy $[B_k,Y_f]=-[Y_f,B_k]\in[\mathfrak l,\mathfrak l]$
oraz $A_k\in\mathfrak l$. Zamieniając kolejność
dwóch czynników pod śladem, możemy zastosować
warunek śladowy. Każdy składnik jest zerowy, więc
$\operatorname{tr}(XY_f)=0$.
Z drugiej strony na $V_i$ operator $XY_f$
ma ślad $m_i\lambda_i f(\lambda_i)$, ponieważ
$N_i$ ma ślad zero (jego wartości własne są zerowe).
Dla każdego $f$ otrzymujemy zatem równość w $K$:
\[
 \sum_{i=1}^s m_i\lambda_i f(\lambda_i)=0.
\]
Lewa strona należy do $E$, bo $f(\lambda_i)\in\mathbb Q$.
Stosując do niej $f$ i korzystając z liniowości
nad $\mathbb Q$, otrzymujemy w $\mathbb Q$
\[
 0=f\Bigl(\sum_i m_i f(\lambda_i)\lambda_i\Bigr)
   =\sum_i m_i f(\lambda_i)^2.
\]
Nie używamy żadnego porządku na $K$, lecz zwykłego
porządku na $\mathbb Q$. Wszystkie składniki
są nieujemnymi liczbami wymiernymi, a $m_i>0$,
więc $f(\lambda_i)=0$ dla każdego $i$.
Gdyby pewne $\lambda_i\ne0$, można by uzupełnić
je do bazy $E$ nad $\mathbb Q$ i wybrać funkcjonał
przyjmujący na nim wartość $1$, co jest sprzeczne
z ostatnią równością. Zatem wszystkie $\lambda_i=0$,
czyli $X$ na każdej uogólnionej przestrzeni własnej
jest równy nilpotentnemu $N_i$. Ponieważ $N_i^d=0$
i przestrzenie te rozpinają $V$, mamy $X^d=0$.

Wykazaliśmy nilpotentność każdego operatora z
$[\mathfrak l,\mathfrak l]$. Wniosek~\ref{aa:cor:engel-flag}
stosuje się do tej podalgebry na $V$ i daje wspólną
postać ściśle górnotrójkątną. Jest ona podalgebrą
algebry nilpotentnej $\mathfrak n_d$
(Przykład~\ref{aa:prz:triangular-lie}), zatem sama
jest nilpotentna, a więc rozwiązalna na mocy
Lematu~\ref{aa:lem:lie-series}. Ponieważ iloraz
$\mathfrak l/[\mathfrak l,\mathfrak l]$ jest abelowy,
Lemat~\ref{aa:lem:solvable-extension} dowodzi
rozwiązalności $\mathfrak l$.
\end{proof}

\begin{twierdzenie}[Kryterium rozwiązalności Cartana]
\label{aa:tw:cartan-solvability}
Niech $\mathfrak g$ będzie skończenie wymiarową algebrą
Liego nad ciałem $K$ charakterystyki zero.
Algebra $\mathfrak g$ jest rozwiązalna wtedy i tylko
wtedy, gdy $\kappa_{\mathfrak g}(x,y)=0$
dla wszystkich $x\in[\mathfrak g,\mathfrak g]$,
$y\in\mathfrak g$.
\end{twierdzenie}
\begin{proof}
Najpierw niech $K$ będzie algebraicznie domknięte.
Stosujemy Twierdzenie~\ref{aa:tw:cartan-linear}
do $\mathfrak l=\operatorname{ad}(\mathfrak g)$
na $V=\mathfrak g$. Ponieważ $\operatorname{ad}$
jest homomorfizmem,
$[\mathfrak l,\mathfrak l]=
\operatorname{ad}([\mathfrak g,\mathfrak g])$.
Warunek śladowy w twierdzeniu jest dokładnie
warunkiem w tezie. Algebra $\mathfrak l$ jest
rozwiązalna wtedy i tylko wtedy, gdy $\mathfrak g$
jest rozwiązalna: w jedną stronę stosujemy
Stwierdzenie~\ref{aa:prop:lie-series-images} do obrazu,
w drugą korzystamy z
$\mathfrak g/Z(\mathfrak g)\simeq\mathfrak l$
i Lematu~\ref{aa:lem:solvable-extension}, ponieważ
centrum jest abelowe.

Dla dowolnego $K$ charakterystyki zero bierzemy
jego domknięcie algebraiczne $L$ i rozszerzamy
skalary: $\mathfrak g_L=\mathfrak g\otimes_K L$,
z nawiasem określonym na tensorach prostych przez
$[x\otimes a,y\otimes b]=[x,y]\otimes ab$.
W bazie $\mathfrak g$ nad $K$ ta sama rodzina
wektorów z tensorem $1$ jest bazą nad $L$.
Nawias jest dobrze określony przez rozszerzenie
dwuliniowe stałych struktury; alternowanie
i tożsamość Jacobiego sprawdzamy na wektorach
bazowych, gdzie są tymi samymi równościami w $K$
odczytanymi w $L$ (alternowanie również wynika
z antysymetrii w charakterystyce zero).
Jeżeli $U,W\subset\mathfrak g$ są podprzestrzeniami,
to ich rozszerzenia skalarów utożsamiamy z
podprzestrzeniami $\mathfrak g_L$: bazę każdej
z nich można uzupełnić do bazy $\mathfrak g$.
Zachodzi równość
$[U\otimes_K L,W\otimes_K L]=[U,W]\otimes_K L$:
nawiasy tensorów prostych dają inkluzję w prawo,
a generatory $[u,w]\otimes a=[u\otimes1,w\otimes a]$
dają odwrotną. Indukcja pokazuje
$(\mathfrak g_L)^{(k)}=\mathfrak g^{(k)}\otimes L$.
Podprzestrzeń po rozszerzeniu skalarów znika
dokładnie wtedy, gdy przed rozszerzeniem jest zerowa,
co wynika z zachowania liczby elementów bazy.
Rozwiązalność jest więc równoważna przed i po
rozszerzeniu.

Macierze $\operatorname{ad}_{x\otimes1}$ w bazie
nad $L$ mają te same wpisy co macierze
$\operatorname{ad}_x$ nad $K$. Forma Killinga
na $\mathfrak g_L$ jest zatem rozszerzeniem
dwuliniowym $\kappa_{\mathfrak g}$.
Dokładniej, dla $x,y\in\mathfrak g$ i $a,b\in L$
mamy $\kappa_{\mathfrak g_L}(x\otimes a,y\otimes b)
=ab\,\kappa_{\mathfrak g}(x,y)$, z prawą stroną
odczytaną w $L$. Jeśli forma znika na
$[\mathfrak g,\mathfrak g]\times\mathfrak g$,
to znika po rozszerzeniu skalarów na
$[\mathfrak g_L,\mathfrak g_L]\times\mathfrak g_L$
przez dwuliniowość. Odwrotnie podstawiamy
$a=b=1$ i korzystamy z iniektywności $K\hookrightarrow L$.
Kryterium nad $L$ oraz wykazana równoważność
rozwiązalności przed i po rozszerzeniu dają tezę nad $K$.
\end{proof}

\section{Kryterium półprostoty Cartana}

\begin{twierdzenie}[Kryterium półprostoty Cartana]
\label{aa:tw:cartan-semisimple}
Niech $\mathfrak g$ będzie skończenie wymiarową algebrą
Liego nad ciałem charakterystyki zero.
Algebra $\mathfrak g$ jest półprosta wtedy i tylko
wtedy, gdy jej forma Killinga jest niezdegenerowana.
\end{twierdzenie}
\begin{proof}
Załóżmy najpierw półprostotę. Niech
$\mathfrak k=\ker\kappa_{\mathfrak g}$.
Z Lematu~\ref{aa:lem:lie-orthogonal-ideals} jest to
ideał. Na mocy Stwierdzenia~\ref{aa:prop:killing-properties}
forma Killinga $\mathfrak k$ jest ograniczeniem
formy Killinga $\mathfrak g$ do $\mathfrak k$,
więc jest identycznie zerowa. Kryterium
rozwiązalności Cartana (Twierdzenie~\ref{aa:tw:cartan-solvability})
stosuje się do $\mathfrak k$ i daje rozwiązalność.
Z półprostoty wynika $\mathfrak k=0$.

Odwrotnie, niech $\kappa_{\mathfrak g}$ będzie
niezdegenerowana. Gdyby radykał rozwiązalny
$\mathfrak r$ był niezerowy, jego ostatni niezerowy
wyraz szeregu pochodnego $\mathfrak a=
\mathfrak r^{(j)}$ byłby niezerowym abelowym
ideałem $\mathfrak g$. Aby uzasadnić własność
ideału w całej $\mathfrak g$, zaczynamy od
$\mathfrak r\triangleleft\mathfrak g$ i stosujemy
indukcyjnie fakt z dowodu Lematu~\ref{aa:lem:lie-series},
że nawias dwóch ideałów jest ideałem.
Dla $x\in\mathfrak a$ mamy
$\operatorname{ad}_x(\mathfrak g)\subset\mathfrak a$
oraz $\operatorname{ad}_x(\mathfrak a)=0$.
Dla dowolnego $y\in\mathfrak g$ operator
$\operatorname{ad}_y$ zachowuje $\mathfrak a$.
Zatem $\operatorname{ad}_x\operatorname{ad}_y$
ma obraz w $\mathfrak a$ i znika na $\mathfrak a$.
W bazie uzupełniającej bazę $\mathfrak a$
jego oba bloki diagonalne są zerowe, a więc
$\kappa_{\mathfrak g}(x,y)=0$.
Stąd $\mathfrak a\subset\ker\kappa_{\mathfrak g}=0$,
sprzeczność. Radykał jest zerowy.
\end{proof}

\section{Rozkład na ideały proste}

\begin{twierdzenie}[Rozkład algebry półprostej]
\label{aa:tw:semisimple-decomposition}
Niech $\mathfrak g$ będzie skończenie wymiarową
półprostą algebrą Liego nad ciałem $K$ charakterystyki
zero. Dla każdego ideału $\mathfrak a$ mamy rozkład
$\mathfrak g=\mathfrak a\oplus\mathfrak a^\perp$
na komutujące ideały, gdzie dopełnienie jest
brane względem formy Killinga.
W szczególności $\mathfrak g$ jest sumą prostą
skończenie wielu ideałów prostych. I odwrotnie,
suma prosta skończenie wielu algebr prostych nad
$K$ jest półprosta.
\end{twierdzenie}
\begin{proof}
Forma Killinga jest niezdegenerowana na mocy
Twierdzenia~\ref{aa:tw:cartan-semisimple}.
Z Lematu~\ref{aa:lem:lie-orthogonal-ideals}
$\mathfrak a^\perp$ jest ideałem oraz
$[\mathfrak a,\mathfrak a^\perp]=0$.
Przecięcie $\mathfrak a\cap\mathfrak a^\perp$
jest ideałem i jest abelowe, bo nawias dwóch
jego elementów należy do
$[\mathfrak a,\mathfrak a^\perp]=0$.
Półprostota wymusza zerowość przecięcia.
Niezdegenerowana forma indukuje izomorfizm
$\mathfrak g\to\mathfrak g^*$, $x\mapsto\kappa(x,\cdot)$.
Ograniczenie funkcjonałów do $\mathfrak a$
jest surjektywne (każdy funkcjonał przedłużamy
przez wybór bazy uzupełniającej).
Złożenie obu odwzorowań ma jądro
$\mathfrak a^\perp$ oraz obraz $\mathfrak a^*$.
Twierdzenie o wymiarze obrazu i jądra daje
$\dim\mathfrak a^\perp=\dim\mathfrak g-\dim\mathfrak a$.
Wraz z zerowym przecięciem dowodzi to rozkładu.
Ograniczenie formy do każdego składnika jest
niezdegenerowane, ponieważ element składnika
ortogonalny do własnego składnika jest również
ortogonalny do drugiego i tym samym do całej
$\mathfrak g$. Formy te są formami Killinga
składników na mocy Stwierdzenia~\ref{aa:prop:killing-properties}.
Kryterium półprostoty daje półprostotę obu składników.

Indukcja po wymiarze: dla wymiaru zero bierzemy
pustą sumę prostą. Niezerowa półprosta algebra
nie może być abelowa. Jeśli nie ma właściwego
niezerowego ideału, jest prosta. W przeciwnym
razie wybieramy taki ideał $\mathfrak a$,
a oba składniki rozkładu mają mniejszy wymiar
i są półproste. Indukcja daje rozkłady na ideały
proste w składnikach. Każdy ideał jednego
składnika jest także ideałem w całej algebrze:
nawias z drugim składnikiem jest zerowy,
a z pierwszym pozostaje w tym ideale.
Złożenie rozkładów daje tezę.

Odwrotnie niech $\mathfrak g=\bigoplus_i\mathfrak g_i$,
gdzie $\mathfrak g_i$ są proste. Dla ideału
rozwiązalnego $\mathfrak r$ rzut na każdy składnik
jest ideałem w tym składniku: dla $x\in\mathfrak g_i$
wybieramy jego przeciwobraz przez rzut, a następnie
stosujemy własność ideału $\mathfrak r$ i
homomorficzność rzutu. Obraz jest rozwiązalny.
Ponieważ każda algebra prosta jest półprosta,
każdy taki obraz jest zerowy. Wszystkie współrzędne
elementu $\mathfrak r$ znikają, więc $\mathfrak r=0$.
\end{proof}

\begin{wniosek}
\label{aa:cor:semisimple-perfect}
Niech $\mathfrak g$ będzie skończenie wymiarową
półprostą algebrą Liego nad ciałem charakterystyki
zero. Wtedy $[\mathfrak g,\mathfrak g]=\mathfrak g$.
\end{wniosek}
\begin{proof}
Z Twierdzenia~\ref{aa:tw:semisimple-decomposition}
mamy sumę prostą komutujących ideałów prostych.
Każdy składnik ma niezerową algebrę pochodną,
bo jest nieabelowy, a ta jest jego ideałem,
więc jest równa składnikowi. Nawias dwóch różnych
składników jest zerowy, a nawiasy wewnętrzne
rozpinają wszystkie składniki. Teza zachodzi
również dla pustej sumy, czyli algebry zerowej.
\end{proof}

\section{Derywacje algebry półprostej}

\begin{twierdzenie}[Derywacje wewnętrzne]
\label{aa:tw:semisimple-derivations}
Niech $\mathfrak g$ będzie skończenie wymiarową
półprostą algebrą Liego nad ciałem charakterystyki
zero. Każda derywacja $D$ algebry $\mathfrak g$
ma postać $\operatorname{ad}_z$ dla dokładnie
jednego $z\in\mathfrak g$.
\end{twierdzenie}
\begin{proof}
Forma Killinga jest niezdegenerowana na mocy
Twierdzenia~\ref{aa:tw:cartan-semisimple}.
Funkcjonał $x\mapsto\operatorname{tr}
(D\operatorname{ad}_x)$ jest liniowy,
zatem istnieje dokładnie jeden $z\in\mathfrak g$
taki, że dla wszystkich $x$
$\operatorname{tr}(D\operatorname{ad}_x)=\kappa(z,x)$:
jest to surjektywność i iniektywność izomorfizmu
z $\mathfrak g$ do $\mathfrak g^*$ wyznaczonego
przez formę niezdegenerowaną w skończonym wymiarze.
Połóżmy $D_0=D-\operatorname{ad}_z$.
Jest derywacją na mocy
Stwierdzenia~\ref{aa:prop:lie-derivations}
i spełnia $\operatorname{tr}(D_0\operatorname{ad}_x)=0$
dla każdego $x$. Dla dowolnych $x,y$ z tej samej
tożsamości o derywacjach mamy
$[D_0,\operatorname{ad}_x]=\operatorname{ad}_{D_0x}$.
Stąd
\[
\begin{aligned}
 \kappa(D_0x,y)
 &=\operatorname{tr}([D_0,\operatorname{ad}_x]
                   \operatorname{ad}_y)\\
 &=\operatorname{tr}(D_0[\operatorname{ad}_x,
                                  \operatorname{ad}_y])\\
 &=\operatorname{tr}(D_0\operatorname{ad}_{[x,y]})=0.
\end{aligned}
\]
Druga równość to cykliczność śladu, trzecia
to własność reprezentacji dołączonej.
Niezdegenerowanie formy daje $D_0x=0$ dla
każdego $x$, więc $D=\operatorname{ad}_z$.
Jednoznaczność wynika także z $Z(\mathfrak g)=0$.
\end{proof}

\section{Formy Killinga klasycznych algebr}

\begin{twierdzenie}
\label{aa:tw:classical-killing}
Niech $K$ będzie ciałem charakterystyki zero,
a $n\ge1$. W oznaczeniach
Przykładu~\ref{aa:prz:classical-lie} zachodzą wzory
\[
\begin{aligned}
 \kappa_{\mathfrak{gl}_n}(X,Y)&=
  2n\operatorname{tr}(XY)-2\operatorname{tr}(X)\operatorname{tr}(Y),\\
 \kappa_{\mathfrak{sl}_n}(X,Y)&=2n\operatorname{tr}(XY),\\
 \kappa_{\mathfrak{so}_n}(X,Y)&=(n-2)\operatorname{tr}(XY),\\
 \kappa_{\mathfrak{sp}_{2n}}(X,Y)&=(2n+2)\operatorname{tr}(XY),
\end{aligned}
\]
gdzie w każdym wierszu $X,Y$ należą do wskazanej
algebry, a ślady po prawej stronie są śladami
w reprezentacjach naturalnych.
\end{twierdzenie}
\begin{proof}
Na $M_n(K)$ piszemy $\operatorname{ad}_X=L_X-R_X$,
gdzie $L_X(A)=XA$, $R_X(A)=AX$.
W bazie jednostek macierzowych mamy
\[
 \operatorname{tr}_{M_n}(L_A)=n\operatorname{tr}(A),
 \quad \operatorname{tr}_{M_n}(R_A)=n\operatorname{tr}(A),
 \quad \operatorname{tr}_{M_n}(L_A R_B)
       =\operatorname{tr}(A)\operatorname{tr}(B).
\]
Istotnie współczynnik przy $E_{ij}$ w $AE_{ij}$
to $a_{ii}$, w $E_{ij}A$ to $a_{jj}$, a w
$AE_{ij}B$ to $a_{ii}b_{jj}$; sumujemy po $i,j$.
Ponieważ $L_XL_Y=L_{XY}$, $R_XR_Y=R_{YX}$
i lewe mnożenia komutują z prawymi,
rozwinięcie $(L_X-R_X)(L_Y-R_Y)$ daje pierwszy wzór.
Algebra $\mathfrak{sl}_n$ jest ideałem w $\mathfrak{gl}_n$,
więc Stwierdzenie~\ref{aa:prop:killing-properties}
daje drugi wzór przez ograniczenie.

Dla pozostałych dwóch przypadków niech $V=K^m$
i niech $b$ będzie niezdegenerowaną formą
dwuliniową symetryczną (gdy $m=n$, macierz $I_n$)
lub alternującą (gdy $m=2n$, macierz $J$).
Niech $\mathfrak l=\{X:b(Xu,v)+b(u,Xv)=0\
\forall u,v\in V\}$. Są to dokładnie rozważane
algebry ortogonalna i symplektyczna.
Odwzorowanie
\[
 T:V\otimes V\to\operatorname{End}_K(V),
 \qquad T(u\otimes v)(w)=b(v,w)u
\]
jest izomorfizmem: niezdegenerowanie $b$ daje
izomorfizm $V\to V^*$, $v\mapsto b(v,\cdot)$,
a w dowolnej bazie $e_i$ i bazie dualnej $e^j$
operatory $w\mapsto e^j(w)e_i$ tworzą bazę
jednostek macierzowych.
Dla $X\in\mathfrak l$ oraz $u,v,w\in V$ mamy
\[
 [X,T(u\otimes v)](w)
 = b(v,w)Xu-b(v,Xw)u
 = b(v,w)Xu+b(Xv,w)u.
\]
Zatem $T$ przeplata działanie dołączone z działaniem
$S_X=X\otimes I+I\otimes X$.

Niech $\tau(u\otimes v)=v\otimes u$.
Podprzestrzenie $V_+$ i $V_-$ odpowiadają wartościom
własnym $+1$ i $-1$ operatora $\tau$.
Ponieważ $\tau^2=I$ i $2$ jest odwracalne,
$P_\pm=(I\pm\tau)/2$ są rzutami na te przestrzenie,
a ich suma to $I$ i iloczyn zero.
Jeśli $b$ jest symetryczna, to
$T(V_-)=\mathfrak l$: dla $A=T(u\otimes v-v\otimes u)$
mamy
\[
 b(Aw,z)+b(w,Az)=
 b(v,w)b(u,z)-b(u,w)b(v,z)
 +b(v,z)b(w,u)-b(u,z)b(w,v)=0.
\]
Takie tensory rozpinają $V_-$, ponieważ współczynniki
w bazie $e_i\otimes e_j$ spełniają $a_{ij}=-a_{ji}$
i $a_{ii}=0$. Zatem obraz jest zawarty w $\mathfrak l$,
a oba wymiary są równe $m(m-1)/2$, co daje równość.
Jeśli $b$ jest alternująca, to dla
$A=T(u\otimes v+v\otimes u)$ analogiczna suma ma postać
\[
 b(v,w)b(u,z)+b(u,w)b(v,z)
 +b(v,z)b(w,u)+b(u,z)b(w,v)=0
\]
na mocy $b(w,u)=-b(u,w)$. Tensory te rozpinają $V_+$,
włącznie z $e_i\otimes e_i$, ponieważ $2$ jest
odwracalne. Wymiar $V_+$ wynosi $m(m+1)/2$;
dla $m=2n$ jest równy $n(2n+1)$, czyli wymiarowi
$\mathfrak l$ z Przykładu~\ref{aa:prz:classical-lie}.
Stąd $T(V_+)=\mathfrak l$.

Operator $S_X$ komutuje z $\tau$ (zamiana czynników
zamienia dwa składniki $S_X$ miejscami), więc
zachowuje obie przestrzenie własne. W konsekwencji
\[
 \operatorname{tr}_{V_\pm}(S_X S_Y)
 =\tfrac12\left(\operatorname{tr}_{V\otimes V}(S_X S_Y)
          \pm\operatorname{tr}_{V\otimes V}(\tau S_X S_Y)\right).
\]
W istocie lewa strona jest śladem $P_\pm S_X S_Y$
na całej przestrzeni, bo na drugim składniku
ten operator znika.

Dla dowolnych operatorów $A,B$ na $V$ mamy
$\operatorname{tr}(A\otimes B)=\operatorname{tr}(A)
\operatorname{tr}(B)$ oraz
$\operatorname{tr}(\tau(A\otimes B))=\operatorname{tr}(AB)$.
Pierwszą równość dostajemy sumując współczynniki
$a_{ii}b_{jj}$ przy $e_i\otimes e_j$, a drugą
sumując współczynniki $a_{ji}b_{ij}$.
Rozwijamy
\[
 S_XS_Y=XY\otimes I+X\otimes Y+Y\otimes X+I\otimes XY.
\]
Zatem jego ślad na całej przestrzeni wynosi
$2m\operatorname{tr}(XY)+2\operatorname{tr}(X)
\operatorname{tr}(Y)$, a po przemnożeniu przez
$\tau$ wynosi $4\operatorname{tr}(XY)$.
W algebrach ortogonalnej i symplektycznej
$\operatorname{tr}(X)=0$: w pierwszej
$X^{\mathsf T}=-X$, a w drugiej bloki diagonalne
mają postać $P,-P^{\mathsf T}$.
Stąd ślad ograniczenia na $V_-$ to
$(m-2)\operatorname{tr}(XY)$, a na $V_+$ to
$(m+2)\operatorname{tr}(XY)$.
Izomorfizm $T$ przeplata te ograniczenia
z reprezentacją dołączoną na $\mathfrak l$,
więc ślady są formami Killinga. Podstawiając
$m=n$ dla $V_-$ i $m=2n$ dla $V_+$,
otrzymujemy ostatnie dwa wzory.
\end{proof}

\begin{przyklad}[Forma Killinga i półprostota modeli klasycznych]
\label{aa:prz:killing-classical-consequences}
Nad ciałem charakterystyki zero forma
$(X,Y)\mapsto\operatorname{tr}(XY)$ na $M_n(K)$
jest niezdegenerowana: $\operatorname{tr}(XE_{ji})=x_{ij}$,
więc znikanie wszystkich takich śladów wymusza $X=0$.
Dla $X\in\mathfrak{sl}_n$ ortogonalnego do całej
$\mathfrak{sl}_n$ testy $Y=E_{ji}$, $i\ne j$,
dają zerowe wpisy poza przekątną, a testy
$Y=E_{ii}-E_{jj}$ dają równe wpisy diagonalne.
Zatem $X=cI_n$ i $nc=0$, czyli $c=0$.
Forma Killinga $\mathfrak{sl}_n$ jest więc
niezdegenerowana dla każdego $n\ge1$.
Dla $\mathfrak{gl}_n$ wzór z
Twierdzenia~\ref{aa:tw:classical-killing}
pokazuje, że $KI_n$ jest w jądrze, a każdy
element jądra jest skalarny na mocy identycznych
testów z macierzami śladu zero. Ponadto
$\mathfrak{gl}_n=KI_n\oplus\mathfrak{sl}_n$
(jawny rozkład $X=(\operatorname{tr}(X)/n)I_n+
(X-(\operatorname{tr}(X)/n)I_n)$). Pierwszy
składnik jest centralnym ideałem, drugi ideałem
półprostym na mocy kryterium Cartana.
Każdy ideał rozwiązalny ma zerowy obraz przez
rzut na $\mathfrak{sl}_n$, więc
$\operatorname{rad}\mathfrak{gl}_n=KI_n$.

Dla $\mathfrak{so}_n$ baza składa się z
$F_{ij}=E_{ij}-E_{ji}$, $i<j$. Ze wzoru
na mnożenie jednostek macierzowych
$\operatorname{tr}(F_{ij}F_{kl})=-2\delta_{ik}\delta_{jl}$
dla uporządkowanych par $i<j,k<l$. Forma śladu
ma diagonalną macierz z niezerowymi wpisami,
a forma Killinga jest niezdegenerowana dla $n\ge3$.
Dla $n=2$ algebra jest jednowymiarowa abelowa,
więc forma Killinga jest zerowa i algebra nie
jest półprosta; dla $n=1$ algebra jest zerowa.

Dla $\mathfrak{sp}_{2n}$ piszemy
$X=\begin{pmatrix}P&Q\\R&-P^{\mathsf T}\end{pmatrix}$,
$Y=\begin{pmatrix}A&B\\C&-A^{\mathsf T}\end{pmatrix}$,
gdzie $Q,R,B,C$ są symetryczne. Mnożenie bloków
i cykliczność śladu dają
\[
 \operatorname{tr}(XY)=2\operatorname{tr}(PA)
                    +\operatorname{tr}(QC)
                    +\operatorname{tr}(RB).
\]
Wybierając kolejno $B=C=0$ i dowolne $A$,
$A=B=0$ i dowolne symetryczne $C$ oraz $A=C=0$
i dowolne symetryczne $B$, wykazujemy, że
ortogonalność do całej algebry wymusza $P=Q=R=0$:
dla $P$ stosujemy test jednostkami macierzowymi,
dla macierzy symetrycznych testy $E_{ii}$
oraz $E_{ij}+E_{ji}$ dają wpisy diagonalne
i dwukrotności wpisów pozadiagonalnych.
Zatem forma śladu jest niezdegenerowana.
Współczynnik $2n+2$ jest niezerowy w charakterystyce
zero, więc forma Killinga także jest niezdegenerowana.
Kryterium~\ref{aa:tw:cartan-semisimple} daje
półprostotę $\mathfrak{sl}_n$, $\mathfrak{so}_n$
dla $n\ge3$ i $\mathfrak{sp}_{2n}$.
Nie dowodzi to jeszcze prostoty wszystkich
niezerowych modeli.
\end{przyklad}

\begin{przyklad}[Macierz formy Killinga na $\mathfrak{sl}_2$]
\label{aa:prz:killing-sl2}
W bazie $(h,e,f)$ z Przykładu~\ref{aa:prz:heisenberg-sl2}
korzystamy ze wzoru $4\operatorname{tr}(XY)$.
Mamy $h^2=I_2$, $ef=E_{11}$, $fe=E_{22}$,
$e^2=f^2=0$, a $he,hf$ mają ślad zero.
Macierz formy i jej wyznacznik wynoszą
\[
 \begin{pmatrix}8&0&0\\0&0&4\\0&4&0\end{pmatrix},
 \qquad -128.
\]
Forma jest niezdegenerowana, choć $\kappa(e,e)=0$
dla $e\ne0$. Niezdegenerowanie nie jest warunkiem
dodatniej określoności.
\end{przyklad}

Uzyskaliśmy rozkład półprostych algebr Liego i
charakterystykę poprzez formę Killinga. Związek
tych algebr z układami korzeni z rozdziału~\ref{aa:ch:root-systems}
wymaga jeszcze rozkładu na przestrzenie korzeniowe,
którym zajmiemy się w dalszej części tomu.
